\subsection{The multiplier rate extracted from UGF}

For $m>0$ define the vertex multiplier
\[
 r_N(m)
 :=\frac{\sqrt{mS_N(m)}}{\omega_N(m)}
 =\sqrt{2m\,g_N(m)}
\quad (m\le N/2),
\qquad
 r_N(m):=0\quad(m>N/2).
\tag{5.5}\label{eq:vertexmultiplier}
\]
The creation part of the finite vertex exponential contains
$r_N(m)a_{-m}/m$; the annihilation part is its adjoint.  Theorem
\ref{thm:UGF} gives $0\le r_N(m)\le\pi/2$.

\begin{lemma}[Uniform quadratic multiplier error]\label{lem:multiplier}
For all even $N\ge16$ and $1\le m\le N/2$,
\[
 |r_N(m)-1|
 \le A_0\frac{m^2}{N^2},
\qquad
 A_0:=4096\left(1+\frac\pi2\right).
\tag{5.6}\label{eq:rNrate}
\]
Consequently, for $1\le m\le N$,
\[
 |r_{2N}(m)-r_N(m)|
 \le A_{\rm TEL}\frac{m^2}{N^2},
\qquad
 A_{\rm TEL}:=\frac54A_0
 =5120\left(1+\frac\pi2\right).
\tag{5.7}\label{eq:dyadicmultiplier}
\]
\end{lemma}

\begin{proof}
Put $t=m/N$.  For $t\le1/64$, every momentum entering
\eqref{eq:structurefactor} lies in $[-\pi/32,\pi/32]$.  Differentiating
the explicit $M=1$ transfer recurrence gives
\[
 \sup_{|p|\le\pi/32}\norm{u'(p)}\le4.
\tag{5.8}\label{eq:edgederivative}
\]
For completeness, the decaying transfer root is
$\rho(p)=1-\cos p$; on this interval
$|\rho|\le1-\cos(\pi/32)$ and $|\rho'|\le\sin(\pi/32)$.
Differentiating the normalized geometric vector gives a norm below
$1$, while the opposite-boundary correction is bounded by
$2|\rho|^{N_y}/(1-|\rho|)$ and its derivative by the corresponding
geometric derivative; their sum is below $4$ for $N_y=8$.
Therefore
\[
 0\le1-\frac{S_N(m)}m
 \le\norm{u(p_{n+m})-u(p_n)}^2
 \le64\pi^2t^2.
\tag{5.9}\label{eq:Srate}
\]
Also, elementary Taylor remainder and
Lemma~\ref{lem:sine} give
\[
 0\le\frac{\pi t}{\sin(\pi t)}-1\le\pi^2t^2
 \qquad(0\le t\le1/2).
\tag{5.10}\label{eq:symbolrate}
\]
Since
$r_N=(\pi t/\sin\pi t)\sqrt{S_N/m}$, these estimates prove
\eqref{eq:rNrate} for $t\le1/64$ with a constant much smaller than
$A_0$.  For $t\ge1/64$, the UGF bound gives
$|r_N-1|\le1+\pi/2\le A_0t^2$, proving the global statement.

For $m\le N/2$, apply \eqref{eq:rNrate} at $N$ and $2N$ and use the
triangle inequality, giving the factor $1+1/4=5/4$.  For
$N/2<m\le N$, $r_N(m)=0$ and
$r_{2N}(m)\le\pi/2\le2\pi m^2/N^2<A_{\rm TEL}m^2/N^2$.
\end{proof}

\begin{lemma}[Lipschitz Wick estimate]\label{lem:WickLipschitz}
Let $r,s$ be two cutoff multiplier sequences bounded by $\pi/2$ and
suppose $|r(m)-s(m)|\le a m^2$.  For the norm-two charge
$\lambda$ and every $f\in C^\infty(S^1)$,
\[
 \norm{\bigl(Y_r(e^\lambda,f)-Y_s(e^\lambda,f)\bigr)
        (1+L_0)^{-3}}
 \le C_\lambda\,a\,\norm{f}_{H^4},
\tag{5.11}\label{eq:WickLipschitz}
\]
where $C_\lambda$ is independent of both cutoffs.  The same estimate
holds for the adjoints.
\end{lemma}

\begin{proof}
Interpolate $r_\tau=(1-\tau)r+\tau s$ and differentiate the finite
normal-ordered exponential.  Every differentiated monomial has exactly
one insertion
\[
 J_{s-r}(z)
 =\sum_{m>0}\frac{s(m)-r(m)}m
   \bigl(a_{-m}z^m-a_mz^{-m}\bigr)
\]
and otherwise the vertex exponential with multiplier $r_\tau$.
The oscillator inequalities
$\norm{a_{\pm m}(1+L_0)^{-1/2}}\le\sqrt{m+1}$, followed by
Cauchy--Schwarz in $m$ and four integrations by parts in the smearing
variable, give
\[
 \norm{:\!J_{s-r}Y_{r_\tau}(e^\lambda)\!:(f)
       (1+L_0)^{-3}}
 \le C_\lambda a\norm{f}_{H^4}.
\]
The constant is uniform because $0\le r_\tau(m)\le\pi/2$, so all
undifferentiated Wick contractions are dominated by
Theorem~\ref{thm:UGF}.  Integrating over $\tau\in[0,1]$ proves
\eqref{eq:WickLipschitz}.  The finite-cutoff calculation is an
absolutely convergent sum after the displayed graph weight; monotone
cutoff removal gives the asserted bound.  Replacing $\lambda$ by
$-\lambda$ and conjugating the coefficients proves the adjoint
statement.
\end{proof}

\begin{theorem}[TEL-a: dyadic charged-field graph defect]
\label{thm:TELa}
The explicit isometries \eqref{eq:Fockisometry} satisfy, for every
$f\in C^\infty(S^1)$,
\[
\begin{aligned}
 &\norm{\bigl(\Psi_{\lambda,2N}(f)V_{N,2N}
       -V_{N,2N}\Psi_{\lambda,N}(f)\bigr)(1+L_{0,N})^{-3}}\\
 &\quad+
 \norm{\bigl(\Psi_{\lambda,2N}(\overline f)^*V_{N,2N}
       -V_{N,2N}\Psi_{\lambda,N}(\overline f)^*\bigr)
       (1+L_{0,N})^{-3}}\\
 &\hspace{35mm}\le
 2C_\lambda A_{\rm TEL}\norm{f}_{H^4}N^{-2}.
\end{aligned}
\tag{5.12}\label{eq:TELa}
\]
Thus TEL-a holds with $p=3$ and $\epsilon=1$.
\end{theorem}

\begin{proof}
After applying $V_{N,2N}$, the coarse field is the fine Fock-space
vertex field with multiplier $r_N$ extended by zero on the new modes.
Apply Lemma~\ref{lem:WickLipschitz} with
$a=A_{\rm TEL}N^{-2}$ from Lemma~\ref{lem:multiplier}; apply it again
to the adjoint.  This proves \eqref{eq:TELa}.
\end{proof}

\begin{lemma}[Summable graph defects imply strong convergence]
\label{lem:telescoping}
The dyadic estimate \eqref{eq:TELa} implies that the transported fields
and adjoints converge on $\Dfin$ along every dyadic cofinal sequence.
\end{lemma}

\begin{proof}
Let $\xi$ have energy at most $E$ and represent it at scale $N_0$.
For $K>L$, insert the compatible isometries and telescope:
\[
\norm{\Psi_{\lambda,2^KN_0}^{\rm tr}(f)\xi
      -\Psi_{\lambda,2^LN_0}^{\rm tr}(f)\xi}
\le
2C_\lambda A_{\rm TEL}\norm{f}_{H^4}(1+E)^3
\sum_{j=L}^{K-1}(2^jN_0)^{-2}.
\tag{5.13}\label{eq:telestimate}
\]
The geometric tail tends to zero.  The adjoint half of
\eqref{eq:TELa} gives graph convergence, and L1 makes the limit
closable.  A dyadic subsystem is cofinal in the scaling parameter, so
this is sufficient for the edge-mode scaling limit.
\end{proof}

\paragraph{Finite audit (numerical support, not proof).}
The companion probe
\path{experiments/tfpt-discovery/mmst_tel_probe.py}
checks \eqref{eq:modeisometry}--\eqref{eq:dyadicmultiplier} at
$N=16,24,32,48,64,96$.  The Gram and compatibility errors are exactly
zero, the measured coefficient envelope is
$\max N^2|r_{2N}(m)-r_N(m)|/m^2=3.420384012310$, and a smooth-test Wick
seminorm fits $N^{-2.022861}$.  A wrong-rate $1/N$ mutant and a
non-isometric embedding mutant are both rejected.

\begin{remark}[Local support]
The Shannon embedding \eqref{eq:modeisometry} is exact in momentum and
is not sharply support-preserving in position.  Theorem~\ref{thm:TELa}
therefore proves graph convergence of the edge-projected fields, not by
itself convergence of sharply localized microscopic algebras.  The
support/locality statement remains explicitly in \textup{(ALG-EXH)}.
\end{remark}

\subsection{TEL-b: the two-edge full-CAR remainder}

The exact Hardy model gives the quarter-twist coefficient
\[
 \norm{P_0-P_{1/4}}_{\HS}^2
 =\frac{2}{\pi^2}\log N+O(1),
\qquad \frac2{\pi^2}=0.202642367\ldots .
\tag{5.14}\label{eq:hardylog}
\]
For the frozen QWZ exhaustion, the computed decomposition is
\[
 C_N^{(0)}-C_N^{(1)}
 =D_{N,\mathrm{top}}^{(1/4)}
  \oplus D_{N,\mathrm{bottom}}^{(1/4)}+R_N,
\tag{5.15}\label{eq:tworemainder}
\]
with
\[
 \max_{N\in\{16,24,32,48,64,96\}}\norm{R_N}_{\HS}
 =2.105079358<2.11.
\tag{5.16}\label{eq:measured211}
\]
A $1/N$ fit predicts $2.107565497$, with maximum fit residual
$2.349\times10^{-5}$.  These values support boundedness of $R_N$ but do
not prove either
\[
 \sup_N\norm{R_N}_{\HS}\le2.11
\tag{5.17}\label{eq:allN211}
\]
for the full-CAR remainder.

\begin{proposition}[Exact Hankel reduction of TEL-b]
\label{prop:HankelReduction}
Let
\[
 C(p)=\id_{(-\infty,0)}(h_{\rm strip}(p))
       +\frac12\id_{\{0\}}(h_{\rm strip}(p))
\]
be the explicit $16\times16$ frozen QWZ Fermi symbol and let
$P_{\rm t},P_{\rm b}$ be the rank-one boundary projections at $p=0$.
The symbol left after the constant-profile Hardy subtraction is
\[
 B(p)=C(p)
 -\id_{(0,\pi)}(p)P_{\rm t}
 -\id_{(-\pi,0)}(p)P_{\rm b}.
\tag{5.18}\label{eq:regularizedsymbol}
\]
Writing $B(p)=\sum_{k\in\Z}B_ke^{ikp}$, every finite Hankel corner
obeys
\[
 \norm{H_N(B)}_{\HS}^2
 \le\sum_{k\ne0}|k|\norm{B_k}_{\mathrm F}^2.
\tag{5.19}\label{eq:PellerHS}
\]
Hence an all-$N$ TEL-b proof by this subtraction would follow from
\[
 16\sum_{k\ne0}|k|\norm{B_k}_{\mathrm F}^2\le C_R^2.
\tag{5.20}\label{eq:Hankeltarget}
\]
\end{proposition}

\begin{proof}
In the untwisted Hardy basis, multiplication by $B$ has matrix entries
$B_{j-k}$.  Its off-diagonal Hardy corner has entries $B_{j+k+1}$.
Counting the pairs $(j,k)\in\mathbb N_0^2$ with $j+k+1=\ell$ gives the exact
Hilbert--Schmidt identity
\[
 \norm{H(B)}_{\HS}^2
 =\sum_{\ell\ge1}\ell\norm{B_\ell}_{\mathrm F}^2.
\]
Finite corners are compressions and cannot increase this norm.  The
opposite edge uses the negative coefficients.  There are at most four
corner/orientation contributions; the triangle inequality gives the
safe factor sixteen after squaring in \eqref{eq:Hankeltarget}.  This is
the $p=2$ case of the Peller $B^{1/2}_{2,2}$ criterion \cite{Peller}.
\end{proof}

\begin{remark}[Second-angle verdict: \texttt{TEL-B-EXTERNAL}]
The new Hankel route does not certify the requested constant.  On a
$4096$-point symbol grid the companion probe finds
\[
 \sum_{k\ne0}|k|\norm{B_k}_{\mathrm F}^2
 =4.482956087,\qquad
 4\norm{B}_{\dot H^{1/2}}=8.469196974>3.
\]
The constant edge-profile subtraction cancels the infrared jump at
$p=0$ but exposes a Brillouin-seam ultraviolet contribution; the raw
symbol seminorm on the same grid is $2.570232705$.  Thus neither a
$C_R=2.5$ nor a $C_R=3$ proof follows, and the value is numerical rather
than an all-mode enclosure.  UGF cannot repair this because it controls
only the scalar edge branch, whereas \eqref{eq:regularizedsymbol} is a
matrix-valued transverse cancellation.  This was the second distinct
in-house attempt; by the stop rule TEL-b is now typed
\texttt{TEL-B-EXTERNAL}.  No all-$N$ constant is claimed.
\end{remark}

\begin{remark}[Round-3 supersession, 2026-09-05]
The preceding verdict records the historical raw-Hankel attempt.  It is not the
current TEL-b residue.  Native $256$-bit Arb/Acb arithmetic on $5812$ rational
cells now proves
\[
 \mathrm{TV}_p(F'')\le15.8176959629889243576\ldots<15.928674,
 \qquad K_3^{\mathrm{TV}}<2.517464500835,
\]
including the exact Hilbert--Schmidt cut jump $2$
(\texttt{v1022\_telb\_round3\_certificates.py}).  Separately, a $100$-bit
certificate over $3888$ analytic rectangles proves
$C_\star<0.1525$, and the Peano/Euler argument proves
$\sup_k|\varphi_N(k)|<0.509/N$ for every even $N\ge16$
(\texttt{v1025\_telb\_c1\_c2a\_certificates.py}).  These close the TV, C1 and
C2a inputs. Round 4 additionally proves CF $<0.065697$, DG $<0.042959$,
and the complete relaxed norm $\norm{R_N}_{\HS}<2.995906<3$ for every even
$N\ge16$, including the finite-grid aliases
(\texttt{v1026\_telb\_round4\_closure.py}). The remaining microscopic
algebra identification is not a consequence of this norm theorem.
\end{remark}

\subsection{Third angle: certified finite part plus a dyadic tail}

There are two natural full-CAR one-particle embeddings from circumference
$N$ to $2N$.  The low-energy embedding $J_N^{\rm low}$ preserves the
integer Fourier label $n$, while the physical-momentum dilation
$J_N^{\rm dil}$ sends $n\mapsto2n$.  The proposed tail mechanism would
require, for one of them,
\[
 \norm{R_{2N}-J_NR_NJ_N^*}_{\HS}\le DN^{-2}
\tag{5.21}\label{eq:Rtailproposed}
\]
with one explicit $D$.  The reverse triangle inequality would then give
\[
 \left|\norm{R_{2N}}_{\HS}-\norm{R_N}_{\HS}\right|
 \le DN^{-2},
\tag{5.22}\label{eq:Rnormtail}
\]
and the dyadic tail from $N_0$ would be at most
$DN_0^{-2}/(1-1/4)$.

\begin{proposition}[The TEL-a embedding does not supply the proposed
remainder tail]\label{prop:tailfailure}
On the frozen finite matrices, neither natural embedding exhibits the
defect required by \eqref{eq:Rtailproposed}.  The measured values are
\[
\begin{array}{c|c|c}
N\to2N&
\norm{R_{2N}-J_N^{\rm low}R_N(J_N^{\rm low})^*}_{\HS}&
\norm{R_{2N}-J_N^{\rm dil}R_N(J_N^{\rm dil})^*}_{\HS}\\ \hline
16\to32&2.071346866&2.066044074\\
32\to64&2.067361892&2.072036965\\
64\to128&2.065611659&2.075227063\\
96\to192&2.065062404&2.076335188
\end{array}
\tag{5.23}\label{eq:Rdefecttable}
\]
Thus the certified-tail argument cannot inherit the proved TEL-a
$N^{-2}$ comparison.
\end{proposition}

\begin{proof}
In the untwisted Fourier basis, construct $R_N$ exactly as in
\eqref{eq:tworemainder}; both embeddings are coordinate isometries.
For either $J$,
\[
 \norm{R_{2N}-JR_NJ^*}_{\HS}^2
 =\norm{R_{2N}}_{\HS}^2+\norm{R_N}_{\HS}^2
 -2\Re\ip{J^*R_{2N}J}{R_N}_{\HS}.
\]
The companion probe evaluates this identity without forming $J$ and
reproduces the established finite norms.  The table is numerical
evidence, not an interval theorem; it shows that the proposed small
TEL-a-type defect is absent on the tested scales.  Therefore
high-precision recalculation of the same finite values cannot certify
the missing infinite tail.
\end{proof}

\begin{remark}[The strictly smaller dyadic residue]
The scalar norm increments are instead
\[
 0.007486497094,\quad0.003790483818,\quad
 0.001913049365,\quad0.001279839965
\]
for the four rows of \eqref{eq:Rdefecttable}; the products $N$ times
these increments are
\[
 0.119783953,\quad0.121295482,\quad
 0.122435159,\quad0.122864637.
\]
This suggests, but does not prove,
\[
 \left|\norm{R_{2N}}_{\HS}-\norm{R_N}_{\HS}\right|
 \le\frac{A_R}{N}\qquad(N\ge96).
\tag{5.24}\label{eq:Rscalarresidue}
\]
Unlike an adjacent-size $1/N$ estimate, \eqref{eq:Rscalarresidue} is
summable on a dyadic tower.  If proved, it gives
\[
 \sup_{k\ge0}\norm{R_{96\cdot2^k}}_{\HS}
 \le\norm{R_{96}}_{\HS}+\frac{2A_R}{96}.
\tag{5.25}\label{eq:Rdyadicbound}
\]
The edge-mode scaling theorem uses a cofinal dyadic subsystem, so
\eqref{eq:Rdyadicbound}, not a supremum over every integer $N$, is the
minimal required statement.  The six reported finite values are
ordinary double-precision diagnostics.  Arbitrary-precision
\texttt{mpmath} output alone is not a rigorous interval certificate;
because the prerequisite tail theorem failed, no floating-point value
is relabelled as certified.
\end{remark}

\begin{remark}[The dyadic-increment route is retired]
Equations~\eqref{eq:Rscalarresidue}--\eqref{eq:Rdyadicbound} are retained as
an audit trail, not as the current proof route.  The cover/sawtooth reduction
and the round-3 C1/C2a/TV certificates supersede this proposed scalar
increment.  No all-$N$ conclusion is imported from the measurements above.
\end{remark}

\paragraph{Third-angle audit.}
\path{experiments/tfpt-discovery/mmst_telb_tail_probe.py} reproduces all
six established finite norms, evaluates both embeddings through
$N=192$, verifies the reverse-triangle inequality, and rejects an
$N^{-1}$ mutant as an $N^{-2}$ tail.  Its verdict is
\texttt{TAIL-TELESCOPING-FAILS}; all diagnostic and honesty guards pass.

